% ============================================================
%  Informe de avance Laboratorio 4 - Control PID (ball and beam)
%  Compilar con: pdflatex informe_lab4.tex  (dos veces)
% ============================================================
\documentclass[11pt,a4paper]{article}

\usepackage[utf8]{inputenc}
\usepackage[T1]{fontenc}
\usepackage[spanish,es-noquoting]{babel}
\usepackage{lmodern}
\usepackage[margin=2.2cm,top=1.9cm,bottom=1.9cm]{geometry}
\usepackage{xcolor}
\usepackage{enumitem}
\usepackage{listings}
\usepackage{titlesec}
\usepackage{float}
\usepackage{tikz}
\usetikzlibrary{shapes.geometric,arrows.meta,positioning,calc}
\usepackage[colorlinks=true,urlcolor=blue!55!black]{hyperref}

\setlist[itemize]{topsep=2pt,itemsep=1.5pt,parsep=0pt,leftmargin=1.2em}
\setlist[enumerate]{topsep=2pt,itemsep=1.5pt,parsep=0pt,leftmargin=1.4em}
\setlength{\parskip}{3pt}
\pagestyle{empty}

\titleformat{\section}{\normalsize\bfseries}{}{0pt}{}
\titlespacing*{\section}{0pt}{9pt}{2pt}

% --- estilo de los listados de código ---
\lstset{
  language=C,
  basicstyle=\ttfamily\scriptsize,
  commentstyle=\color{black!55}\itshape,
  keywordstyle=\bfseries,
  showstringspaces=false,
  tabsize=2,
  frame=leftline,
  framerule=0.6pt,
  rulecolor=\color{black!40},
  xleftmargin=8pt,
  aboveskip=4pt,
  belowskip=2pt,
  literate={á}{{\'a}}1 {é}{{\'e}}1 {í}{{\'i}}1 {ó}{{\'o}}1 {ú}{{\'u}}1
           {ñ}{{\~n}}1 {¿}{{?`}}1 {→}{{$\to$}}1
}

% --- estilo del diagrama ---
\tikzset{
  paso/.style={rectangle, rounded corners=3pt, draw=black!60, fill=blue!6,
               text width=10.2cm, align=left, inner sep=3pt, font=\footnotesize},
  decision/.style={diamond, aspect=3.5, draw=black!60, fill=orange!12,
                   align=center, inner sep=0pt, font=\footnotesize},
  mecanica/.style={rectangle, rounded corners=3pt, draw=black!60, fill=red!8,
                   text width=3.8cm, align=left, inner sep=3pt, font=\footnotesize},
  flecha/.style={-{Stealth[length=2.2mm]}, thick, black!70}
}

\begin{document}

% ------------------------------------------------------------
\begin{center}
  {\large\bfseries Control PID de un balancín (ball and beam)}\\[2pt]
  {\small Laboratorio 4: informe de avance \textbar{} Sistemas Electrónicos de Instrumentación}\\[2pt]
  {\footnotesize Pedro Villar \quad\textbar\quad 23 de noviembre de 2026}
\end{center}
\vspace{-4pt}
\hrule
\vspace{4pt}

El objetivo es mantener una pelota de ping pong en el centro de una viga. Un HC-SR04 mide dónde
está la pelota, un servo inclina la viga y un PID calcula cada 50\,ms cuánto inclinarla. Todo
corre en una \textbf{STM32F4Discovery}. El PID sigue la implementación de la guía de la
cátedra.

% ------------------------------------------------------------
\section*{El lazo principal}

\begin{lstlisting}
while (1) {
  uint32_t tiempo_actual = HAL_GetTick();

  if ((tiempo_actual - tiempo_anterior) >= TS_MS) {      // Ts = 50 ms
    tiempo_anterior = tiempo_actual;

    float posicion_actual = LeerSensorDistancia();       // cm respecto del centro

    if (modo == 1U && fallas_seguidas < FALLAS_PARA_CENTRAR) {
      pid_bola.banda = banda_cm;                      // editable en vivo
      salida_angulo = PID_Compute(&pid_bola, setpoint_cm,
                                  posicion_actual);
    } else {
      salida_angulo = 0.0f;                           // midiendo o sin pelota:
      PID_Reset(&pid_bola, posicion_actual);          // viga plana y PID en cero
    }
    ActualizarServo(salida_angulo);                 // pide el ángulo, sin saltos
  }

  Servo_Update(tiempo_actual);  // cada 20 ms acerca el servo (150 grados/s)
  Update_Leds();                // rojo: sin pelota, verde: centrada
}
\end{lstlisting}

Es el mismo esquema no bloqueante de la guía: el \texttt{if} del tiempo decide cuándo corre el
lazo y nada frena al resto. \texttt{LeerSensorDistancia()} toma el eco de la medición anterior,
dispara la siguiente y devuelve la mediana de las últimas 5 lecturas, porque el sensor da picos
falsos. El servo no salta al ángulo que pide el PID: se mueve como mucho 150 grados por segundo,
así la viga no da golpes. Las variables \texttt{modo}, \texttt{setpoint\_cm}, \texttt{banda\_cm}
y las ganancias se cambian con la placa andando, desde el depurador.

% ------------------------------------------------------------
\section*{Callbacks de interrupción}

\begin{lstlisting}
// TIM3 en Input Capture: llega un flanco del eco del HC-SR04
void HAL_TIM_IC_CaptureCallback(TIM_HandleTypeDef *htim) {
  hcsr04_tim_ic_callback(&distance_sensor, htim);   // la librería mide el ancho del eco
}

// TIM4 cada 50 ms: manda una muestra por UART para graficar en la PC
void HAL_TIM_PeriodElapsedCallback(TIM_HandleTypeDef *htim) {
  if (htim->Instance == TIM4) {
    // tiempo, posición, setpoint y salida del PID, separados por comas
    int largo = snprintf(linea_grafico, sizeof(linea_grafico),
                         "%lu,%.2f,%.2f,%.2f\r\n", (unsigned long)HAL_GetTick(),
                         posicion_cm, setpoint_cm, salida_angulo);
    HAL_UART_Transmit_IT(&huart2, (uint8_t *)linea_grafico, (uint16_t)largo);
  }
}
\end{lstlisting}

El primero es el del sensor, igual que en el laboratorio 3. El segundo lo agregué para
calibrar: cada 50\,ms manda por USART2 una línea con el tiempo, la posición, el setpoint y la
salida del PID a un adaptador USB-UART. En la PC un script de Python la grafica en tiempo real y
la guarda en un CSV. Usa \texttt{Transmit\_IT}, así la placa no espera a que salga la línea y
el lazo no se atrasa.

% ------------------------------------------------------------
\section*{Lo que me pasó}

Primero calibré la maqueta. El centro de la viga quedó a 33\,cm del sensor y la viga queda
horizontal con un pulso de 1315\,µs. Ahí apareció el primer dato raro: entre 1240 y 1390\,µs la
pelota no se mueve para ningún lado. Son unos 13 grados de servo que no hacen nada.

Después probé el PID y me encontré con tres problemas:

\begin{itemize}
  \item \textbf{El servo temblaba.} Hacía cambios chiquitos todo el tiempo aunque la pelota
        estuviera quieta. El sensor tiene un ruido de unos 0,4\,cm y la derivada lo divide por
        Ts, así que lo agranda. Lo resolví suavizando la velocidad (70\,\% del valor anterior
        y 30\,\% del nuevo), y ese problema quedó resuelto.
  \item \textbf{La pelota se traba.} Cerca del medio recién empieza a rodar con unos 19 grados
        de servo. A 37\,cm del sensor se quedó quieta más de un minuto con el servo al tope
        (35 grados). El brazo del servo inclina muy poco la viga, y por eso tuve que subir el
        límite de 15 a 35 grados.
  \item \textbf{Cuando se destraba, se pasa.} La pelota junta fuerza hasta que vence el roce,
        sale de golpe y cruza el centro con velocidad. Llega al medio, pero tarda muchísimo en
        quedarse quieta.
\end{itemize}

Las ganancias que mejor anduvieron fueron Kp 2,5, Ki 0,3 y Kd 0,8, con una banda de 1\,cm
alrededor del centro. Las saqué a prueba y error, cambiándolas en vivo. No pude aplicar
Ziegler-Nichols como pide la guía: con tanto roce la pelota nunca oscila de
forma pareja, se traba y sale a los saltos. Así que no hay una ganancia crítica para medir.

% ------------------------------------------------------------
\section*{Cómo decidí resolverlo}

Me di cuenta de que estaba ajustando a ojo algo que no había medido. Además, todo
apunta a que el problema es más mecánico que de sintonía: un PID lineal no puede compensar del
todo una zona en la que el servo se mueve y la pelota no. Por eso cambié el orden: primero
medir la planta y después elegir las ganancias.

Para eso dejé el firmware preparado. Simplifiqué el PID (derivada sobre la medida,
anti-windup y nada más), agregué un setpoint que se puede cambiar en vivo y amplié lo que
manda la UART. También armé un script de MATLAB que toma el log de un ensayo de escalón y
hace los pasos del diagrama.

\begin{figure}[H]
\centering
\begin{tikzpicture}[node distance=3mm]
  \node[paso] (p1) {\textbf{1. Ensayo de escalón en la maqueta.}
    Ki = 0 y banda = 0. Con la pelota quieta en el centro, el setpoint pasa de 0 a 8\,cm.
    La UART guarda el log en un CSV.};
  \node[paso, below=of p1] (p2) {\textbf{2. Primera estimación (método de Chen).}
    Del sobrepico y el tiempo de pico sale la ganancia K de la planta: cuánto acelera la
    pelota por cada grado de servo.};
  \node[paso, below=of p2] (p3) {\textbf{3. Ajuste por simulación.}
    MATLAB simula el lazo con la misma lógica del firmware (mediana, retardo, derivada
    suavizada, servo lento) y busca el K y la zona muerta que mejor copian el log.};
  \node[decision, below=4mm of p3] (d) {¿Zona muerta grande?};
  \node[mecanica, anchor=west] (m) at ($(d -| p3.east) + (6mm,0)$) {\textbf{Compensarla en el software.}
    A la salida del PID se le suman los grados de zona muerta, hacia el lado que tiene
    que ir la pelota.};
  \node[paso, below=8mm of d] (p4) {\textbf{4. Probar ganancias sobre el modelo.}
    Se simulan varios Kp y Kd y se eligen los que no oscilan y casi no se pasan.};
  \node[paso, below=of p4] (p5) {\textbf{5. Validar en la maqueta.}
    Otro escalón con las ganancias nuevas. La respuesta real tiene que parecerse a la
    simulada. Recién ahí se sube Ki de a poco y se graban las ganancias.};

  \draw[flecha] (p1) -- (p2);
  \draw[flecha] (p2) -- (p3);
  \draw[flecha] (p3) -- (d);
  \draw[flecha] (d) -- node[above, font=\small]{sí} (m);
  \draw[flecha] (d) -- node[right, font=\small]{no} (p4);
  \draw[flecha] (p4) -- (p5);
  \draw[flecha] (m.south) |- (p4.east);
\end{tikzpicture}
\end{figure}

La fórmula de Chen supone un sistema de segundo orden ideal. Mi lazo tiene filtros,
retardos y un servo con velocidad limitada, así que ese número solo sirve como punto de partida.
El que cuenta es el del paso 3, porque simula el lazo tal como está programado. Además, ese paso
mide la zona muerta, que es justamente lo que sospecho que me limita.

% ------------------------------------------------------------
\section*{Análisis y conclusiones}

Hasta ahora la parte electrónica y de software funciona: la medición no bloquea el lazo, el
filtro saca los picos del sensor, el servo se mueve sin golpes y puedo ver y guardar todo en
la PC mientras la placa corre. Lo que todavía no logré es que la pelota se quede quieta en el
centro de forma estable. La llevo al medio, pero tarda mucho y no siempre se queda.

La conclusión principal es que el límite no está en el código sino en la maqueta. Hay unos 13
grados de servo en los que la pelota no se mueve y hacen falta casi 19 para que arranque. Con
eso, cualquier ganancia queda entre trabarse o pasarse. Aprendí que antes de sintonizar un PID
conviene medir la planta: probar valores a ojo me llevó mucho tiempo sin una razón clara de por
qué no andaba. La maqueta es fija y no la puedo modificar, así que si la zona muerta es grande
la tengo que compensar en el software. El próximo paso es hacer el ensayo de escalón, ver cuánto
da la zona muerta y elegir las ganancias con el modelo, para llegar al final con el lazo
funcionando.

\vspace{6pt}
\noindent El script de MATLAB del análisis y el código completo del proyecto (firmware, librería
del PID y gráfico en Python) están en el repositorio
\href{https://github.com/PedroMVillar/Sist.Instrumentacion/tree/lab4/laboratorio-4-control-PID}{\texttt{github.com/PedroMVillar/Sist.Instrumentacion}},
en la rama \texttt{lab4}, carpeta \texttt{laboratorio-4-control-PID}.

\end{document}
